% Capillarità: (a) acqua, menisco concavo, risale di più nel tubo più sottile;
% (b) mercurio, menisco convesso, si abbassa; (c) angolo di contatto.
\begin{tikzpicture}[>=Stealth, font=\small]
  % (a) acqua
  \fill[formula!20] (0,0) rectangle (3,0.8);
  \draw[very thick] (0,1.2) -- (0,0) -- (3,0) -- (3,1.2);
  % tubo largo
  \fill[formula!20] (0.6,0.3) -- (0.6,1.55) .. controls (0.75,1.38) and (1.05,1.38) .. (1.2,1.55) -- (1.2,0.3) -- cycle;
  \draw[thick] (0.6,0.3) -- (0.6,3.2) (1.2,0.3) -- (1.2,3.2);
  \draw[formula, thick] (0.6,1.55) .. controls (0.75,1.38) and (1.05,1.38) .. (1.2,1.55);
  % tubo sottile
  \fill[formula!20] (1.9,0.3) -- (1.9,2.65) .. controls (1.97,2.55) and (2.13,2.55) .. (2.2,2.65) -- (2.2,0.3) -- cycle;
  \draw[thick] (1.9,0.3) -- (1.9,3.2) (2.2,0.3) -- (2.2,3.2);
  \draw[formula, thick] (1.9,2.65) .. controls (1.97,2.55) and (2.13,2.55) .. (2.2,2.65);
  \draw[formula, thick] (0,0.8) -- (0.6,0.8) (1.2,0.8) -- (1.9,0.8) (2.2,0.8) -- (3,0.8);
  \draw[dashed, richiamo] (2.2,0.8) -- (2.8,0.8);
  \draw[dashed, richiamo] (2.2,2.58) -- (2.8,2.58);
  \draw[|<->|] (2.65,0.8) -- (2.65,2.58) node[midway, right] {$h$};
  \node at (1.5,-0.45) {(a) acqua};

  % (b) mercurio
  \begin{scope}[xshift=4.2cm]
    \fill[richiamo!45] (0,0) rectangle (2.2,1.4);
    \draw[very thick] (0,1.8) -- (0,0) -- (2.2,0) -- (2.2,1.8);
    \fill[white] (0.8,0.3) rectangle (1.4,1.4);
    \fill[richiamo!45] (0.8,0.3) -- (0.8,0.85) .. controls (0.95,1.05) and (1.25,1.05) .. (1.4,0.85) -- (1.4,0.3) -- cycle;
    \draw[thick] (0.8,0.3) -- (0.8,2.6) (1.4,0.3) -- (1.4,2.6);
    \draw[richiamo, thick] (0.8,0.85) .. controls (0.95,1.05) and (1.25,1.05) .. (1.4,0.85);
    \draw[dashed, richiamo] (1.4,1.4) -- (2.0,1.4);
    \node at (1.1,-0.45) {(b) mercurio};
  \end{scope}

  % (c) angolo di contatto: calotte sferiche su una superficie piana
  \begin{scope}[xshift=7.2cm, yshift=0.2cm, scale=1.45]
    % goccia che bagna: centro (1,1.1), raggio 0,8, theta = 51 gradi
    \fill[formula!20] (0.3755,1.6) arc[start angle=141.3, end angle=38.7, radius=0.8] -- cycle;
    \draw[formula, thick] (0.3755,1.6) arc[start angle=141.3, end angle=38.7, radius=0.8];
    \draw[thick] (0,1.6) -- (2.2,1.6);
    \draw (0.3755,1.6) -- ++(51.3:0.6);
    \draw (0.6755,1.6) arc[start angle=0, end angle=51.3, radius=0.3];
    \node[font=\footnotesize] at (0.78,1.78) {$\theta$};
    \node[font=\footnotesize, align=center] at (1.0,1.15) {bagna: $\theta < \ang{90}$};
    % goccia che non bagna: centro (1,0,25), raggio 0,45, theta = 124 gradi
    \fill[richiamo!45] (0.626,0) arc[start angle=213.8, end angle=-33.8, radius=0.45] -- cycle;
    \draw[richiamo, thick] (0.626,0) arc[start angle=213.8, end angle=-33.8, radius=0.45];
    \draw[thick] (0,0) -- (2.2,0);
    \draw (0.626,0) -- ++(123.8:0.5);
    \draw (0.876,0) arc[start angle=0, end angle=123.8, radius=0.25];
    \node[font=\footnotesize] at (0.62,0.42) {$\theta$};
    \node[font=\footnotesize, align=center] at (1.0,-0.35) {non bagna: $\theta > \ang{90}$};
    \node at (1.0,-0.95) {(c)};
  \end{scope}
\end{tikzpicture}
